% CameraModule
\section{【中級】CameraModule — PSRAMバッファ + フレーム管理}

\begin{patternbox}{CameraModule で学べるパターン}
\begin{itemize}
    \item PSRAMへの大容量バッファ配置
    \item フレームバッファのライフサイクル管理（取得→使用→解放）
    \item Data構造体にメタデータのみ格納し、生データはモジュール内に保持
    \item 多数のピンをConfig構造体で一元管理
\end{itemize}
\end{patternbox}

\subsection{Configの大量ピン管理}

\begin{lstlisting}[style=cppstyle, caption={CameraModule.h — Config（多数のピン）}]
struct CameraConfig {
    int8_t pwdnPin, resetPin;           // 電源/リセット
    int8_t xclkPin, siodPin, siocPin;   // クロック/SCCB
    int8_t d0Pin, d1Pin, d2Pin, d3Pin;  // DVPデータバス
    int8_t d4Pin, d5Pin, d6Pin, d7Pin;
    int8_t vsyncPin, hrefPin, pclkPin;  // 同期
    int     xclkFreqHz;
    uint8_t frameSize, pixFormat;
    int     jpegQuality;
    size_t  fbCount;
};
\end{lstlisting}

\begin{pointbox}[Config構造体の威力]
カメラモジュールはピンが16本以上ある。
コンストラクタの引数で16個のピンを渡すと順序を間違える危険があるが、
Config構造体なら\texttt{.d0Pin = 11}のように名前で指定できる。
Config構造体のピンアサインが\texttt{ProjectConfig.h}に集約されるため、
回路設計者はこのファイルだけでピンの確認・変更ができる。
\end{pointbox}

\subsection{フレームバッファの管理}

\begin{lstlisting}[style=cppstyle, caption={CameraModule — フレームバッファ管理}]
// updateInput()はメタデータだけをSystemDataに書き込む
void CameraModule::updateInput(SystemData& data) {
    // 保持中のフレームがあれば先に解放
    releaseFrame();

    camera_fb_t* fb = esp_camera_fb_get();
    if (fb) {
        data.camera.frameReady = true;
        data.camera.width = fb->width;
        data.camera.height = fb->height;
        data.camera.frameSize = fb->len;
        _frameBuffer = fb;  // void*で保持（ヘッダ依存回避）
        _frameHeld = true;
    }
}

// ロジックフェーズから生データにアクセスする場合
const uint8_t* CameraModule::getFrameBuffer() const {
    if (!_frameHeld || !_frameBuffer) return nullptr;
    return ((camera_fb_t*)_frameBuffer)->buf;
}
void CameraModule::releaseFrame() {
    if (_frameHeld && _frameBuffer) {
        esp_camera_fb_return((camera_fb_t*)_frameBuffer);
        _frameBuffer = nullptr;
        _frameHeld = false;
    }
}
\end{lstlisting}

\begin{pointbox}[Data構造体に生データを入れない]
\texttt{camera\_fb\_t*}（フレームバッファポインタ）はSystemDataに格納しない。
SystemDataに入れると、他モジュールがフレームの解放タイミングを制御できてしまう。
代わりに\texttt{getFrameBuffer()} / \texttt{releaseFrame()}でアクセスを制御し、
\texttt{CameraData}にはメタデータ（サイズ、解像度等）のみを格納する。
\end{pointbox}

\begin{notebox}[PSRAMの自動配置]
\texttt{fbCount >= 2}の場合、esp\_cameraライブラリが自動的にPSRAMにバッファを配置する。
Config構造体の\texttt{fbCount}で制御でき、モジュール内部でPSRAM判定のコードを書く必要はない。
\end{notebox}

% ─── API仕様 ───────────────────────────────
\subsection{API仕様}


\begin{apibox}{CameraModule --- OV5640カメラモジュール}
\textbf{定義ファイル:} \texttt{lib/CameraModule/CameraModule.h / .cpp}\\
\textbf{分類:} 入力モジュール（\texttt{updateInput()}のみオーバーライド）\\
\textbf{対象デバイス:} OV5640（ESP32-S3 DVP接続）

ESP32-S3のDVPインターフェース経由でOV5640カメラからフレームを取得するモジュール。
\texttt{esp\_camera}ライブラリを使用し、フレームバッファの取得・解放を管理する。
フレームの生データは\texttt{SystemData}には格納せず、専用メソッドでアクセスする。
\end{apibox}

\subsubsection{機能概要}

\begin{itemize}
    \item DVP 16ピンインターフェースによるカメラ接続
    \item JPEG/RAWフォーマット対応（\texttt{pixFormat}で選択）
    \item フレームバッファの取得・解放API（\texttt{getFrameBuffer()} / \texttt{releaseFrame()}）
    \item フレームメタデータ（解像度・サイズ）の\texttt{SystemData}への書き込み
    \item \texttt{deinit()}によるカメラリソース解放
\end{itemize}

% --- Config ---
\subsubsection{Config構造体}

\begin{funcblock}{CameraConfig}
DVPピンアサインとカメラ動作パラメータを定義する。

\begin{longtable}{l l l p{4.2cm}}
\toprule
\textbf{メンバ} & \textbf{型} & \textbf{制約・既定値} & \textbf{説明} \\
\midrule
\endfirsthead
\toprule
\textbf{メンバ} & \textbf{型} & \textbf{制約・既定値} & \textbf{説明} \\
\midrule
\endhead
\midrule
\multicolumn{4}{r}{\small 次ページに続く} \\
\endfoot
\bottomrule
\endlastfoot
\multicolumn{4}{l}{\textbf{--- DVPピンアサイン ---}} \\
\texttt{pwdnPin}    & \texttt{int8\_t} & -1で未使用         & POWER DOWNピン \\
\texttt{resetPin}   & \texttt{int8\_t} & -1で未使用         & ハードウェアRESETピン \\
\texttt{xclkPin}    & \texttt{int8\_t} & ---                & XCLKクロック出力ピン \\
\texttt{siodPin}    & \texttt{int8\_t} & ---                & SCCB SDA（カメラ設定用I2C） \\
\texttt{siocPin}    & \texttt{int8\_t} & ---                & SCCB SCL \\
\texttt{d0Pin}--\texttt{d7Pin} & \texttt{int8\_t} & --- & データバス D0--D7（8ピン） \\
\texttt{vsyncPin}   & \texttt{int8\_t} & ---                & 垂直同期ピン \\
\texttt{hrefPin}    & \texttt{int8\_t} & ---                & 水平参照ピン \\
\texttt{pclkPin}    & \texttt{int8\_t} & ---                & ピクセルクロックピン \\
\midrule
\multicolumn{4}{l}{\textbf{--- カメラ動作設定 ---}} \\
\texttt{xclkFreqHz}  & \texttt{int}     & 通常: 20000000   & XCLKクロック周波数 [Hz] \\
\texttt{frameSize}   & \texttt{uint8\_t} & 例: 5 (QVGA)    & フレームサイズ（\texttt{framesize\_t}の値） \\
\texttt{pixFormat}   & \texttt{uint8\_t} & 例: 4 (JPEG)    & ピクセルフォーマット（\texttt{pixformat\_t}の値） \\
\texttt{jpegQuality} & \texttt{int}     & 0--63（低い=高画質）& JPEG圧縮品質 \\
\texttt{fbCount}     & \texttt{size\_t} & 推奨: 2          & フレームバッファ数 \\
\end{longtable}
\end{funcblock}

% --- Data ---
\subsubsection{Data構造体}

\begin{funcblock}{CameraData}
フレームのメタデータを保持する。フレームの生データ（ピクセルバッファ）は含まない。

\begin{longtable}{l l l p{4.8cm}}
\toprule
\textbf{メンバ} & \textbf{型} & \textbf{初期値} & \textbf{説明} \\
\midrule
\endhead
\texttt{isValid}    & \texttt{bool}     & \texttt{false} & カメラ初期化が成功していれば\texttt{true} \\
\texttt{frameReady} & \texttt{bool}     & \texttt{false} & 新しいフレームが取得済みであれば\texttt{true} \\
\texttt{width}      & \texttt{uint16\_t} & \texttt{0}    & フレーム幅 [px] \\
\texttt{height}     & \texttt{uint16\_t} & \texttt{0}    & フレーム高さ [px] \\
\texttt{frameSize}  & \texttt{size\_t}  & \texttt{0}     & フレームデータサイズ [bytes] \\
\bottomrule
\end{longtable}
\end{funcblock}

\begin{cautionbox}[フレームデータのアクセス方法]
フレームの生データ（ピクセルバッファ）は大容量のためSystemDataには格納しない。
ロジックフェーズで\texttt{getFrameBuffer()}を呼び出してポインタを取得し、
使用後は必ず\texttt{releaseFrame()}で解放すること。
解放しないとフレームバッファが枯渇し、次のフレーム取得がブロックされる。
\end{cautionbox}

% --- メソッド ---
\subsubsection{メソッド}

\begin{funcblock}{CameraModule(const CameraConfig\& config)}
\begin{tabularx}{\textwidth}{l X}
\toprule
\textbf{項目} & \textbf{内容} \\
\midrule
引数   & \texttt{config} --- \texttt{CameraConfig}へのconst参照 \\
動作   & Configを\texttt{\_config}にコピー保持する \\
\bottomrule
\end{tabularx}
\end{funcblock}

\begin{funcblock}{bool init()}
\begin{tabularx}{\textwidth}{l X}
\toprule
\textbf{項目} & \textbf{内容} \\
\midrule
戻り値     & \texttt{bool} --- カメラ初期化成功で\texttt{true}、失敗で\texttt{false} \\
動作       & Configからesp\_cameraの設定構造体を構築し、\texttt{esp\_camera\_init()}を呼び出す \\
失敗条件   & カメラ未接続、ピン設定誤り、PSRAMメモリ不足等 \\
\bottomrule
\end{tabularx}
\end{funcblock}

\begin{funcblock}{void updateInput(SystemData\& data)}
\begin{tabularx}{\textwidth}{l X}
\toprule
\textbf{項目} & \textbf{内容} \\
\midrule
書き込み先     & \texttt{data.camera} \\
タイミング制御 & 毎ループ実行 \\
動作           & 前回のフレームが保持中であれば解放し、新しいフレームを\texttt{esp\_camera\_fb\_get()}で取得。メタデータ（幅、高さ、サイズ）を\texttt{data.camera}に書き込む \\
\bottomrule
\end{tabularx}
\end{funcblock}

\begin{funcblock}{const uint8\_t* getFrameBuffer() const}
\begin{tabularx}{\textwidth}{l X}
\toprule
\textbf{項目} & \textbf{内容} \\
\midrule
戻り値   & フレームデータへのポインタ。フレーム未取得時は\texttt{nullptr} \\
用途     & ロジックフェーズでフレームデータにアクセスする（BLE送信、SD保存、画像処理等） \\
有効期間 & 次の\texttt{updateInput()}呼び出しまで、または\texttt{releaseFrame()}呼び出しまで \\
\bottomrule
\end{tabularx}
\end{funcblock}

\begin{funcblock}{void releaseFrame()}
\begin{tabularx}{\textwidth}{l X}
\toprule
\textbf{項目} & \textbf{内容} \\
\midrule
動作   & 保持中のフレームバッファを\texttt{esp\_camera\_fb\_return()}で解放する。フレーム未保持時は何もしない \\
\bottomrule
\end{tabularx}
\end{funcblock}

\begin{funcblock}{void deinit()}
\begin{tabularx}{\textwidth}{l X}
\toprule
\textbf{項目} & \textbf{内容} \\
\midrule
動作   & 保持中のフレームを解放し、\texttt{esp\_camera\_deinit()}でカメラリソースを解放する \\
\bottomrule
\end{tabularx}
\end{funcblock}

\begin{cautionbox}[PSRAMの必要性]
\texttt{fbCount}を2以上に設定する場合、PSRAMが必要である。
PSRAMがない環境では\texttt{fbCount = 1}に制限し、フレームサイズも小さくする必要がある。
\texttt{platformio.ini}で\texttt{board\_build.psram = enabled}を設定すること。
\end{cautionbox}

% --- 使用例 ---
\subsubsection{使用例}

\begin{lstlisting}[style=cppstyle, caption={CameraModuleの使用例}]
// ProjectConfig.h
const CameraConfig CAMERA_CONFIG = {
    .pwdnPin = -1, .resetPin = -1, .xclkPin = 15,
    .siodPin = 4,  .siocPin = 5,
    .d0Pin = 11, .d1Pin = 9, .d2Pin = 8, .d3Pin = 10,
    .d4Pin = 12, .d5Pin = 18, .d6Pin = 17, .d7Pin = 16,
    .vsyncPin = 6, .hrefPin = 7, .pclkPin = 13,
    .xclkFreqHz  = 20000000,
    .frameSize   = 5,   // QVGA (320x240)
    .pixFormat   = 4,   // JPEG
    .jpegQuality = 12,
    .fbCount     = 2,   // ダブルバッファ（PSRAM必須）
};

// ロジックフェーズでの使用
void applyPattern(SystemData& data) {
    if (data.camera.frameReady) {
        const uint8_t* buf = camera.getFrameBuffer();
        size_t len = data.camera.frameSize;
        // BLE送信やSD保存などの処理
        camera.releaseFrame();
    }
}
\end{lstlisting}
